% ----------------------
\subsection{Section 14 (June 1829)}
% ----------------------
	\label{DC14}
	
	% -----
	\subsubsection{Note on the JSP and BC versions}
	% -----
	
		The version featured in the JSP comes from BC, so for this section they are the same; the BC is therefore not included in the alternate versions below. This section doesn't appear at all in RB1.
	
	% -----
	\subsubsection{Historical Background}
	% -----
		\label{DC14-HB}
		
		% --
		\paragraph{JSP}
		% --
		
			``JS dictated this revelation for David Whitmer after Whitmer traveled to Harmony, Pennsylvania, and helped JS and Oliver Cowdery move to Peter Whitmer Sr.'s house in Fayette, New York. The move was undertaken to facilitate completion of the Book of Mormon translation, which had caused JS trouble with both Isaac Hale and other residents in Harmony. David Whitmer had heard about JS in the fall of 1828 when rumors about JS's retrieval of the plates circulated widely in the Palmyra, New York, region. While in Palmyra, Whitmer discussed the story of JS and the plates with Cowdery and others but initially supposed it was a rumor. Cowdery, who was acquainted with the Joseph Smith Sr. family, told Whitmer that ``there must be some truth in the story of the plates, and that he intended to investigate the matter.'' Several months later, Cowdery informed Whitmer that he intended to go to Harmony to see JS about the plates. On his journey, Cowdery stopped in Fayette to visit Whitmer and promised to write to him regarding what he learned. Cowdery sent at least three letters to Whitmer from Harmony, the last of which directed Whitmer to bring his wagon to Harmony and move JS and Cowdery to Fayette.
			
			``Aided by what they considered divine manifestations, the Whitmers quickly became followers of JS. David Whitmer later recounted that during their journey to Fayette, he, Cowdery, and JS briefly encountered a ``pleasant, nice looking old man'' whom JS identified by revelation as a heavenly messenger transporting the plates. Whitmer also recalled that soon after their arrival in Fayette, his mother, Mary Musselman Whitmer, was met ``by the same old man,'' who showed her the plates.%
				\footnote{%
					% \label{}%
					The sources for this seem pretty old, compared to the actual events; here is the JSP's note: ``Joseph F. Smith, New York City, NY, to John Taylor et al., [Salt Lake City, Utah Territory], 17 Sept. 1878, draft, Joseph F. Smith, Papers, CHL; Stevenson, Journal, 23 Dec. 1877, 9 Feb. 1886, and 2 Jan. 1887. While the Joseph F. Smith account did not identify the person transporting the plates, Stevenson's accounts variously identified him as `one of the Nephites' and `one of the 3 Nephites.' (See also [Andrew Jenson], `Eight Witnesses,' \textit{Historical Record}, Oct. 1888, 621.).''
					}
			JS's history recorded that ``David, John, and Peter Whitmer Jr became our zealous friends and assistants in the work; And being anxious to know their respective duties, and having desired with much earnestness that I should enquire of the Lord concerning them, I did so, through the means of the Urim and Thummin and obtained for them in succession the folowing Revelations,'' referring to this revelation and the June 1829 revelations directed to John Whitmer and Peter Whitmer Jr. The five revelations dated June 1829 are presented herein in the order found in the index to Revelation Book 1.''%
				\footnote{%
					% \label{}%
					So the order is \hyperref[DC14]{DC 14}, \hyperref[DC18]{DC 18}, \hyperref[DC15]{DC 15}, \hyperref[DC16]{DC 16}, and \hyperref[DC17]{DC 17}.
					}
			
	% -----
	\subsubsection{Versions}
	% -----
		\label{DC14-V}
	
		\begin{tabular}{p{1in}p{2in}p{1in}p{2in}}
			JSP		& \hyperref[EMS-1828-1829-JS-JUN-1829-A]{June 1829}		& BC & Chapter 12, 32--33 \\
			DC 1835	& Section 39, 169--170									& TS & 3:20 (15 August 1842), 885 \\
			MS		& 3:10 (February 1843), 164								& DC 1844 & Section 39, 251--252 \\
		\end{tabular}
		
		\medskip
		
		See the \href{https://drive.google.com/file/d/1goAvNz8jS4R8slYfMG_db55Ty2z_vmgK/view?usp=sharing}{sections comparison} document for more information about the version history.

	% -----
	\subsubsection{Section 14:1} % *DC14-1 
	% -----
		\label{DC14-1}

		\footnote{%
			% \label{}%
			BC has the following head: ``1 \textit{A Revelation given to David [Whitmer], in Fayette, New-York, June, 1829.}'' DC 1835 has, ``\textit{Revelation given to David Whitmer, June, 1829.}.'' DC 1844 has, ``\textit{Revelation given to David Whitmer, June, 1829.}''
			}%
		A GREAT and marvelous work is about to come forth unto the children of men.

		\begin{framed}
		\scriptsize
			\begin{compactdesc}
				\item[JSP] A GREAT and marvelous work is about to come forth unto the children of men: 
				\item[DC 1835] 1 A great and marvelous work is about to come forth unto the children of men: 
				\item[DC 1844] 1 A great and marvellous work is about to come forth unto the children of men: 
			\end{compactdesc}
		\end{framed}

	% -----
	\subsubsection{Section 14:2} % *DC14-2 
	% -----
		\label{DC14-2}

			\footnote{%
				% \label{}%
				I won't go point out the similarities again, but clearly this section draws a lot from other sections given shortly before this one in 1829.
				}%
		Behold, I am God; give heed to my word, which is quick and powerful, sharper than a two-edged sword, to the dividing asunder of both joints and marrow; therefore give heed unto my word.

		\begin{framed}
		\scriptsize
			\begin{compactdesc}
				\item[JSP] behold I am God, and give heed to my word, which is quick and powerful, sharper than a two-edged sword, to the dividing asunder of both joints and marrow: therefore, give heed unto my word.
				\item[DC 1835] behold I am God, and give heed to my word, which is quick and powerful, sharper than a two-edged sword, to the dividing asunder of both joints and marrow: therefore, give heed unto my word.
				\item[DC 1844] behold I am God, and give heed to my word, which is quick and powerful, sharper than a two-edged sword, to the dividing asunder of both joints and marrow: therefore give heed unto my word.
			\end{compactdesc}
		\end{framed}

	% -----
	\subsubsection{Section 14:3} % *DC14-3 
	% -----
		\label{DC14-3}

		Behold, the field is white already to harvest; therefore, whoso desireth to reap let him thrust in his sickle with his might, and reap while the day lasts, that he may treasure up for his soul everlasting salvation in the kingdom of God.

		\begin{framed}
		\scriptsize
			\begin{compactdesc}
				\item[JSP] 2 Behold the field is white already to harvest, therefore, whoso desireth to reap, let him thrust in his sickle with his might, and reap while the day lasts, that he may treasure up for his soul everlasting salvation in the kingdom of God: 
				\item[DC 1835] 2 Behold the field is white already to harvest, therefore, whoso desireth to reap, let him thrust in his sickle with his might, and reap while the day lasts, that he may treasure up for his soul everlasting salvation in the kingdom of God: 
				\item[DC 1844] 2 Behold the field is white already to harvest, therefore, whoso desireth to reap, let him thrust in his sickle with his might, and reap while the day lasts, that he may treasure up for his soul, everlasting salvation in the kingdom of God: 
			\end{compactdesc}
		\end{framed}

	% -----
	\subsubsection{Section 14:4} % *DC14-4 
	% -----
		\label{DC14-4}

		Yea, whosoever will thrust in his sickle and reap, the same is called of God.

		\begin{framed}
		\scriptsize
			\begin{compactdesc}
				\item[JSP] Yea, whosoever will thrust in his sickle and reap, the same is called of God: 
				\item[DC 1835] yea, whosoever will thrust in his sickle and reap, the same is called of God: 
				\item[DC 1844] yea, whosoever will thrust in his sickle and reap, the same is called of God:
			\end{compactdesc}
		\end{framed}

	% -----
	\subsubsection{Section 14:5} % *DC14-5 
	% -----
		\label{DC14-5}

		Therefore, if you will ask of me you shall receive; if you will knock it shall be opened unto you.

		\begin{framed}
		\scriptsize
			\begin{compactdesc}
				\item[JSP] therefore, if you will ask of me you shall receive; if you will knock it shall be opened unto you.
				\item[DC 1835] therefore, if you will ask of me you shall receive; if you will knock it shall be opened unto you.
				\item[DC 1844] therefore, if you will ask of me you shall receive; if you will knock it shall be opened unto you.
			\end{compactdesc}
		\end{framed}

	% -----
	\subsubsection{Section 14:6} % *DC14-6 
	% -----
		\label{DC14-6}

		Seek to bring forth and establish my Zion.  Keep my commandments in all things.

		\begin{framed}
		\scriptsize
			\begin{compactdesc}
				\item[JSP] 3 Seek to bring forth and establish my Zion.--- Keep my commandments in all things, 
				\item[DC 1835] 3 Seek to bring forth and establish my Zion. Keep my commandments in all things, 
				\item[DC 1844] 3 Seek to bring forth and establish my Zion. Keep my commandments in all things,
			\end{compactdesc}
		\end{framed}

	% -----
	\subsubsection{Section 14:7} % *DC14-7 
	% -----
		\label{DC14-7}

		And, if you keep my commandments and endure to the end you shall have eternal life,
			\footnote{%
				% \label{}%
				Compare \hyperref[1NEPHI-15-36]{1 Nephi 15:36}, \hyperref[DC6-13]{DC 6:13},
				}%
		which gift is the greatest of all the gifts of God.

		\begin{framed}
		\scriptsize
			\begin{compactdesc}
				\item[JSP] and if you keep my commandments, and endure to the end, you shall have eternal life; which gift is the greatest of all the gifts of God.
				\item[DC 1835] and if you keep my commandments and endure to the end, you shall have eternal life; which gift is the greatest of all the gifts of God.
				\item[DC 1844] and if you keep my commandments and endure to the end, you shall have eternal life; which gift is the greatest of all the gifts of God.
			\end{compactdesc}
		\end{framed}

	% -----
	\subsubsection{Section 14:8} % *DC14-8 
	% -----
		\label{DC14-8}

		And it shall come to pass, that if you shall ask the Father in my name, in faith believing, you shall receive the Holy Ghost, which giveth utterance, that you may stand as a witness of the things of which you shall both hear and see, and also that you may declare repentance unto this generation.

		\begin{framed}
		\scriptsize
			\begin{compactdesc}
				\item[JSP] 4 And it shall come to pass, that if you shall ask the Father in my name, in faith believing, you shall receive the Holy Ghost, which giveth utterance, that you may stand as a witness of the things of which you shall both hear and see; and also, that you may declare repentance unto this generation.
				\item[DC 1835] 4 And it shall come to pass, that if you shall ask the Father in my name, in faith believing, you shall receive the Holy Ghost, which giveth utterance, that you may stand as a witness of the things of which you shall both hear and see; and also, that you may declare repentance unto this generation.
				\item[DC 1844] 4 And it shall come to pass, that if you shall ask the Father in my name, in faith believing, you shall receive the Holy Ghost, which giveth utterance, that you may stand as a witness of the things of which you shall both hear and see; and also that you may declare repentance unto this generation.
			\end{compactdesc}
		\end{framed}

	% -----
	\subsubsection{Section 14:9} % *DC14-9 
	% -----
		\label{DC14-9}

		Behold, I am Jesus Christ, the Son of the living God, who created the heavens and the earth, a light which cannot be hid in darkness;

		\begin{framed}
		\scriptsize
			\begin{compactdesc}
				\item[JSP] 5 Behold I am Jesus Christ the Son of the living God, which created the heavens and the earth; a light which cannot be hid in darkness: 
				\item[DC 1835] 5 Behold I am Jesus Christ the Son of the living God, who created the heavens and the earth; a light which cannot be hid in darkness: 
				\item[DC 1844] 5 Behold I am Jesus Christ the Son of the living God, who created the heavens and the earth; a light which cannot be hid in darkness:
			\end{compactdesc}
		\end{framed}

	% -----
	\subsubsection{Section 14:10} % *DC14-10 
	% -----
		\label{DC14-10}

		Wherefore, I must bring forth the fulness of my gospel from the Gentiles unto the house of Israel.

		\begin{framed}
		\scriptsize
			\begin{compactdesc}
				\item[JSP] wherefore, I must bring forth the fulness of my gospel from the Gentiles unto the house of Israel. 
				\item[DC 1835] wherefore, I must bring forth the fulness of my gospel from the Gentiles unto the house of Israel. 
				\item[DC 1844] wherefore, I must bring forth the fulness of my gospel from the Gentiles unto the house of Israel. 
			\end{compactdesc}
		\end{framed}

	% -----
	\subsubsection{Section 14:11} % *DC14-11 
	% -----
		\label{DC14-11}

		And behold, thou art David, and thou art called to assist; which thing if ye do, and are faithful, ye shall be blessed both spiritually and temporally, and great shall be your reward. Amen.

		\begin{framed}
		\scriptsize
			\begin{compactdesc}
				\item[JSP] And behold thou art David, and thou art called to assist: Which thing if ye do, and are faithful, ye shall be blessed both spiritually and temporally, and great shall be your reward: Amen.
				\item[DC 1835] And behold thou art David, and thou art called to assist: which thing if ye do, and are faithful, ye shall be blessed both spiritually and temporally, and great shall be your reward. Amen.
				\item[DC 1844] And behold thou art David, and thou art called to assist: which thing if ye do, and are faithful, ye shall be blessed both spiritually and temporally, and great shall be your reward: Amen.
			\end{compactdesc}
		\end{framed}